\documentclass[11pt]{article}

\usepackage[margin=1in]{geometry}
\usepackage{amsmath,amssymb,amsthm}
\usepackage{booktabs}
\usepackage{algorithm}
\usepackage{algpseudocode}
\usepackage{hyperref}
\usepackage{xcolor}

\newtheorem{proposition}{Proposition}
\newtheorem{remark}{Remark}
\DeclareMathOperator{\wt}{wt}
\DeclareMathOperator{\rowsp}{rowsp}
\DeclareMathOperator{\Var}{Var}
\newcommand{\F}{\mathbb{F}}
\newcommand{\PL}{P_{\mathrm{L}}}
\newcommand{\HX}{H_X}
\newcommand{\HZ}{H_Z}
\newcommand{\dub}{\hat d_{\mathrm{ub}}}
\newcommand{\deff}{d_{\mathrm{eff}}}
\newcommand{\ceil}[1]{\left\lceil #1 \right\rceil}
\newcommand{\floor}[1]{\left\lfloor #1 \right\rfloor}

\title{\bfseries Estimating the Distance of a Quantum LDPC Code by Decoder
Simulation:\\ the Syndrome-Decoder Upper Bound, with Sub-Threshold Scaling as
Support}
\author{}
\date{\today}

\begin{document}
\maketitle

\begin{abstract}
When a CSS quantum LDPC code is too large or too unstructured to certify its
distance $d$ exactly (Brouwer--Zimmermann and integer programming both become
intractable for dense, high-rate codes), the recognised tool is a decoder-based
\emph{heuristic} that returns an \emph{upper bound}. The standard such method is
the \emph{syndrome-decoder method}: inject random errors, decode, and read off the
weight of the residual logical operator; the minimum over many trials is a
checkable upper bound $\dub\ge d$, tightening monotonically with the number of
trials. We describe this method, its confidence model, and how it is verified by
independent re-simulation. As supporting context we include the sub-threshold
logical-error-rate scaling $\PL\sim(p/p_{\mathrm{th}})^{\ceil{d/2}}$, whose slope
measures the decoder's \emph{effective} distance $\deff$; this is the standard
fault-tolerance suppression law, but $\deff\le d$ and can \emph{under}estimate the
true distance for a sub-optimal decoder, so it characterises performance rather
than finding the code distance. A worked $[[180,20,14]]$ example exhibits all
signals---including a cautionary case where the effective distance reads $12$
while the code distance is $14$.
\end{abstract}

\section{Setup and notation}
\label{sec:setup}

A CSS code on $n$ qubits is given by parity-check matrices
$\HX\in\F_2^{m_X\times n}$, $\HZ\in\F_2^{m_Z\times n}$ with $\HX\HZ^{\mathsf T}=0$.
Write $C_X=\ker\HZ$, $C_Z=\ker\HX$. The nontrivial $X$-logicals are
$C_X\setminus\rowsp\HX$ and
\begin{equation}
  d_X=\min\{\wt(v): v\in C_X\setminus\rowsp\HX\},\qquad
  d_Z=\min\{\wt(v): v\in C_Z\setminus\rowsp\HZ\},
\end{equation}
with $d=\min(d_X,d_Z)$. Fix a $Z$-logical basis $L_Z\in\F_2^{k\times n}$
(computed once from $\HX,\HZ$ by Gaussian elimination); a vector $r\in C_X$ is a
nontrivial $X$-logical iff $L_Z r\neq 0\pmod 2$. We probe $d_X$ with an $X$-error
channel; $d_Z$ is symmetric (swap $X\leftrightarrow Z$). ``The distance'' below
means whichever sector is simulated; report $d=\min(d_X,d_Z)$.

\paragraph{Noise model.} Code-capacity: each qubit independently suffers an $X$
error with probability $p$, perfect syndrome measurement. This isolates the code
from syndrome-extraction-circuit details. (Circuit-level noise measures a
\emph{circuit} effective distance, a different quantity;
see Section~\ref{sec:caveats}.)

\section{The syndrome-decoder distance estimate}
\label{sec:syndrome}

\subsection{Method}
Sample an error $e$, decode it, and inspect the residual. The syndrome is
$\sigma=\HZ e\bmod 2$; a decoder $\mathcal D$ returns a correction $\hat c$ with
$\HZ\hat c=\sigma$, so the residual
\begin{equation}
  r=e\oplus\hat c\ \in\ C_X=\ker\HZ
\end{equation}
\emph{always} lies in the $X$-code. Whenever $L_Z r\neq 0\pmod 2$, the residual is
a genuine nontrivial $X$-logical, and therefore an explicit witness of weight
$\wt(r)\ge d$. Tracking the minimum such weight over many trials gives
\begin{equation}
  \dub=\min_{\text{trials with }L_Z r\neq 0}\wt(r)\ \ge\ d,
  \label{eq:dub}
\end{equation}
a \emph{checkable upper bound} on the distance: each candidate is a concrete
operator one can re-verify ($L_Z r\neq0$ and $\HZ r=0$). This is the
\emph{syndrome-decoder method}, a recognised heuristic distance-finder for large
quantum LDPC codes~\cite{Survey,RoffeBPOSD}; it is the decoder-based cousin of the
random-information-set search of QDistRnd~\cite{QDistRnd}.

\paragraph{Decoder.} qLDPC checks have weight $>2$, so the code is not matchable
and minimum-weight perfect matching does not apply. Use \textbf{BP\,+\,OSD}
(belief propagation with ordered-statistics decoding)~\cite{RoffeBPOSD,PanteleevKalachev};
a reference implementation is the \texttt{ldpc} package. A higher OSD order finds
lower-weight residuals (tighter $\dub$); report the order used.

\paragraph{Error distribution.} Any distribution that drives the decoder into
nontrivial cosets works; two useful choices are (i) code-capacity errors at a rate
near threshold, which produce a steady stream of residual logicals, and (ii)
fixed-weight errors of weight $\approx\ceil{d/2}$, which target the regime where
low-weight logicals appear. The estimate~\eqref{eq:dub} is an upper bound under
any choice.

\begin{algorithm}[t]
\caption{Syndrome-decoder distance upper bound}
\label{alg:syndrome}
\begin{algorithmic}[1]
\Require checks $\HZ$, logical basis $L_Z$, decoder $\mathcal D$, trials $T$
\State $\dub\gets \infty$
\For{$t=1,\dots,T$}
  \State sample an error $e$;\quad $\hat c\gets\mathcal D(\HZ e\bmod 2)$;\quad
         $r\gets e\oplus\hat c$
  \If{$L_Z r\neq0\pmod2$ \textbf{and} $\wt(r)<\dub$} \ $\dub\gets\wt(r)$,\ store $r$
  \EndIf
\EndFor
\State \Return $\dub$ and its witness $r$ \Comment{$d\le\dub$, re-checkable}
\end{algorithmic}
\end{algorithm}

\subsection{Confidence}
\label{sec:syndrome-conf}
$\dub$ is a non-increasing function of the trial count $T$ that is bounded below
by $d$, so it converges from above. Confidence is the \emph{saturation} of this
running minimum: once $\dub$ has not decreased over a large number of additional
trials, low-weight logicals are being missed only with small probability. Two
quantitative handles:
\begin{itemize}
\item \emph{Hitting probability.} If a weight-$w$ logical is reachable as a
residual with per-trial probability $q$, the chance of missing it in $T$ trials is
$(1-q)^T\le e^{-qT}$; observing no improvement over $T$ trials bounds the missed
mass by the rule of three, $q\lesssim 3/T$ at $95\%$.
\item \emph{Cross-method agreement.} Running an independent upper-bound search
(QDistRnd / random information set~\cite{QDistRnd}, or QDistEvol~\cite{Survey})
and obtaining the same value strengthens the claim that $\dub=d$.
\end{itemize}
The method returns an \emph{upper} bound; a matching \emph{lower} bound (hence a
certificate $d=\dub$) requires an exact method---Brouwer--Zimmermann, MaxSAT, or
integer programming~\cite{Survey}---which is feasible only for low-rate or
structured codes (Section~\ref{sec:caveats}).

\section{Supporting context: effective distance from sub-threshold scaling}
\label{sec:slope}

The logical error rate of a distance-$d$ code under a good decoder obeys the
standard fault-tolerance suppression law
\begin{equation}
  \PL(p)\;\sim\;A\,p^{\ceil{d/2}}\quad(p\to0),
  \qquad\Longleftrightarrow\qquad
  \log\PL = \ceil{d/2}\,\log p + \log A + O(p),
  \label{eq:scaling}
\end{equation}
because a decoder that corrects all errors of weight $\le\floor{(d-1)/2}$ can only
fail at weight $\ge\ceil{d/2}$, and the lowest surviving power of $p$ dominates as
$p\to0$. The slope of $\log\PL$ versus $\log p$ thus tends to $\ceil{d/2}$. This
is the law behind exponential error suppression and the Google ``below threshold''
experiments ($\varepsilon_L\propto\Lambda^{-(d+1)/2}$)~\cite{DKLP,Google}.

\subsection{It measures the \emph{effective}, not the code, distance}
The exponent recovered from a fit is set by the \emph{decoder}, not by the code
geometry alone. Define the effective distance $\deff$ by
$\ceil{\deff/2}=$ asymptotic slope. A maximum-likelihood decoder gives $\deff=d$;
a sub-optimal decoder that occasionally fails on an error of weight $<\ceil{d/2}$
introduces a lower power of $p$ and \emph{lowers} the slope, so
\begin{equation}
  \deff\le d .
  \label{eq:effle}
\end{equation}
This gap is documented and can be large: approximate decoders have been reported
at effective distances of only $43$--$65\%$ of the true distance~\cite{Google}.
Consequently the slope characterises \emph{performance} and yields a
\emph{lower-leaning} estimate of $d$; it is not, by itself, a reliable
code-distance finder, and the distance-finding literature does not list it among
distance algorithms~\cite{Survey}.

\subsection{Fitting and its statistics}
With $x_i=\log p_i$, $y_i=\log\widehat\PL(p_i)$ and failure counts $f_i$, the
delta method gives $\Var(y_i)\approx 1/f_i$, so a weighted least-squares fit with
weights $w_i=f_i$ yields slope $\hat\beta$ and
\begin{equation}
  \mathrm{SE}(\hat\beta)=\Bigl[\textstyle\sum_i f_i\,(x_i-\bar x_w)^2\Bigr]^{-1/2},
  \qquad
  \deff\approx 2\hat\beta\ \ (\pm\,2\,\mathrm{SE}(\hat\beta)\ \text{at }1\sigma).
  \label{eq:wls}
\end{equation}
Two practical hazards: (i) finite-$p$ curvature makes the local slope
\emph{rise} toward $\ceil{\deff/2}$ as $p$ decreases, so fit only the saturated
window $p\in[p_{\mathrm{th}}/6,\,p_{\mathrm{th}}/3]$; and (ii) under-resolved
points (few $f_i$) must be dropped---they both inflate $\mathrm{SE}$ and, if
included, bias the fit. To shrink the interval, increase the failures $f_i$
(adaptive sampling to a fixed failure count) and widen the $\log p$ lever arm by
placing points at the extremes of the valid window.

\subsection{The error-weight onset: the sharpest probe of $\deff$ (and a warning)}
\label{sec:onset}
A direct, curvature-free way to read $\ceil{\deff/2}$ is to estimate
$q_w=\Pr(\text{failure}\mid\text{error of weight exactly }w)$ for increasing $w$:
$q_w$ jumps from $0$ to positive at $w^\star=\ceil{\deff/2}$, and a zero-failure
run at $w^\star-1$ gives $q_{w^\star-1}\le 3/N$ ($95\%$). This is sharper and
cheaper than the slope fit, but it keys on the \emph{lowest} weight at which the
decoder \emph{ever} fails and is therefore the most sensitive to
sub-optimality~\eqref{eq:effle}: it returns $\deff$, which may sit well below $d$.
The discriminator is to inspect the residual weights of those onset failures
(Section~\ref{sec:syndrome}): if the decoder's lowest-weight failures still land
on \emph{high}-weight logicals, the low onset is decoder sub-optimality, not a
light logical---see the worked example.

\section{Speed and scale}
\label{sec:speed}

Both methods are Monte Carlo and embarrassingly parallel: independent trials with
distinct seeds, aggregated by a minimum (syndrome decoder) or by summed counts
(slope). Profiling BP\,+\,OSD on the worked-example code (single core): decode-only
$\approx 5.0\times10^4$/s, full sampling loop $\approx 2.6\times10^4$/s---so
roughly half the single-core wall-time is Python overhead around an
already-native (C++) decoder. The two effective levers, with measured speed-ups on
an $11$-core machine:

\paragraph{Multiprocessing (the dominant, near-free lever).} Run independent
worker processes, each with its own decoder and seed; aggregate. For the
syndrome-decoder method this took the throughput from $1.7\times10^4$ to
$1.46\times10^5$ trials/s, an $\mathbf{8.6\times}$ speed-up on $9$ workers, and the
extra trials tightened $\dub$ from $15$ to the correct $14$. One platform
caveat dominates the result: use the \texttt{fork} start method, not
\texttt{spawn}. Under \texttt{spawn} (the macOS default) every worker re-imports
NumPy and the decoder library, which halved the speed-up to $\sim4.5\times$;
\texttt{fork} inherits the parent image and recovers the near-linear scaling.

\paragraph{Threaded native kernel (a C++ port done cheaply).} The
random-information-set upper bound (the QDistRnd-style search) is pure linear
algebra and is already a C++ extension; making it multi-threaded is a small change
with \emph{no decoder-quality risk}. Splitting the trials across \texttt{std::thread}s
per side---over a self-contained, const-input core routine, so no locks---and
releasing the GIL ran at $1.27\times10^5$ trials/s, $\mathbf{7.4\times}$ over one
core, in a single Python call with none of the pickling or process-startup cost of
multiprocessing, and still returns the correct $\dub$. This $7.4\times$ is the
\emph{hardware ceiling}, not a software limit: the test machine has $5$ performance
and $6$ efficiency cores, and the efficiency cores run at roughly a third of the
speed, so the available throughput is $\approx 5 + 6\times0.4 \approx 7.4$
performance-core-equivalents---saturated by using all $11$ threads (fewer leaves
cores idle; more does not help). Two refinements gave little here: hoisting the
per-thread basis recompute (the basis is cheap next to the trials loop) and
load-balancing the static split (the heterogeneous cores, not the split, set the
ceiling). The more compute-bound BP\,+\,OSD multiprocessing reached a slightly
higher $8.6\times$ because its workload draws more fully on the efficiency cores;
the achievable factor is workload-dependent, but ``use all the cores'' is the rule.

\paragraph{What not to do.} Re-implementing BP\,+\,OSD in C++ to remove the
remaining Python overhead is \emph{not} worthwhile: the decoder is already native,
so a rewrite reclaims only the $\sim2\times$ glue factor while risking a weaker OSD
that returns \emph{looser} residual bounds---degrading the very estimate it is
meant to speed up. Prefer multiprocessing around the good decoder, and thread the
native QDistRnd kernel for the random-information-set bound.

\paragraph{Scaling and the rare-event wall.} The syndrome-decoder method needs only
enough trials for $\dub$ to saturate, whereas the slope method's sub-threshold cost
grows as $p^{-\ceil{d/2}}$: even a $20\times$ engine extends the reachable $p$-window
only modestly, a further argument for the upper-bound methods when $d$ is large.

\section{Verification by independent re-simulation}
\label{sec:verify}

A claim of this kind is verified by \emph{re-running the experiment}, not by
re-checking the claimant's arithmetic. Given only the code ($\HX,\HZ$, from which
$L_Z$ is recomputed) plus the reported decoder configuration and seed, an
independent party reconstructs the harness and reproduces the estimate:
for the syndrome-decoder bound, confirm a residual logical of weight $\le\dub$ is
found (and, ideally, re-verify the stored witness directly: $\HZ r=0$,
$L_Z r\neq0$); for the slope, rerun the sweep and confirm the saturated slope. The
submitted counts and fitted numbers are \emph{reproducibility metadata}, not the
object of trust. Re-simulation is minutes-scale, which is acceptable: it is the
honest cost of an independent statistical check, and exact certification, where
feasible at all, is itself minutes-to-hours.

\section{What the method does and does not establish}
\label{sec:caveats}

\begin{itemize}
\item The syndrome-decoder method returns a checkable \emph{upper} bound
$d\le\dub$~\eqref{eq:dub}, with confidence given by the saturation of the running
minimum and by agreement with independent upper-bound searches
(QDistRnd/QDistEvol). It is \emph{not} a certificate: a matching lower bound needs
an exact method (Brouwer--Zimmermann, MaxSAT, integer programming), feasible only
for low-rate or structured codes.
\item The slope (and the error-weight onset) measure the \emph{effective} distance
$\deff\le d$~\eqref{eq:effle}; with a near-ML decoder $\deff=d$, but a sub-optimal
decoder can read substantially lower. Use them for performance characterisation
and as corroboration, not as the primary code-distance estimate.
\item Code-capacity quantities are properties of the \emph{code}. Circuit-level
memory experiments (e.g.\ in Stim) measure a circuit effective distance that folds
in the syndrome-extraction schedule and is generally smaller; use those for
fault-tolerance performance.
\item These simulations are the statistical complement to exact distance
computation. They are most useful exactly where the exact methods fail: dense,
high-rate codes whose minimum-weight-logical ILP has a loose relaxation and whose
information sets are too large to enumerate.
\end{itemize}

\appendix
\section{Worked example: $[[180,20,14]]$}
\label{app:example}
A weight-$8$ two-block (coset) qLDPC code, $X$-sector, BP\,+\,OSD decoder. The
three signals illustrate the hierarchy of this note.

\paragraph{Syndrome-decoder upper bound (primary).} Random errors decoded; the
minimum residual logical weight observed is
\[
  \dub = 14,
\]
in agreement with an independent QDistRnd-style random-information-set search that
returns $14$ and does not improve over $3\times10^5$ trials. Hence $d\le 14$ with
high confidence (two independent upper-bound methods saturate at $14$).

\paragraph{Slope (supporting).} Adaptive sampling to $F\approx 60$ failures,
threshold $p_{\mathrm{th}}\approx 0.085$:
\begin{center}
\begin{tabular}{rrrr}
\toprule
$p$ & $\widehat{\PL}$ & $f/N$ & local slope \\
\midrule
$0.050$ & $5.1\times10^{-2}$ & $102/2{,}000$   & ---  \\
$0.040$ & $1.2\times10^{-2}$ & $73/6{,}000$    & $6.4$ \\
$0.030$ & $1.6\times10^{-3}$ & $60/38{,}000$   & $7.1$ \\
\bottomrule
\end{tabular}
\end{center}
The weighted fit gives slope $6.8\pm0.3$, i.e.\ $\deff\approx 13.6$ (95\% CI
$[12.3,14.8]$)---consistent with $14$ in this window.

\paragraph{Error-weight onset (cautionary).} Fixed-weight sampling, OSD order
$6$--$20$, $10^6$ samples per weight:
\[
  q_5=0\ (\le 3\times10^{-6}),\quad q_6\approx 1\times10^{-5},\quad
  q_7\approx 4\times10^{-5},
\]
so the onset is $w^\star=6$, giving $\deff\in\{11,12\}$---\emph{below} the true
distance. Raising the OSD order from $6$ to $20$ barely moved $q_6$, confirming
this is decoder sub-optimality, not a near-threshold artefact. The discriminator
clinches it: the residual logicals produced by those weight-$6$ failures have
weight $18$--$31$ (minimum $18$), and the lightest residual seen at any onset
weight is $14$. No logical lighter than $14$ is ever exhibited. The onset's
``$12$'' is the effective distance; the code distance is $14$.

\paragraph{Takeaway.} The upper-bound (syndrome-decoder / QDistRnd) value $14$ is
the trustworthy distance estimate; the slope corroborates it; the error-weight
onset measures the BP\,+\,OSD effective distance ($\approx12$) and, taken alone,
would have underestimated $d$. Always cross-check a low onset against the residual
weights.

\section{A family of estimated codes}
\label{app:family}
The primary (syndrome-decoder) method, multiprocessed across all cores and
cross-checked against the threaded QDistRnd bound, applied to a saved family of
weight-$8$ coset codes on the same group ($n=180$). Each $\dub$ is the minimum
residual logical weight over $6\times10^5$ trials per method; the two independent
upper-bound methods agree on every code.

\begin{center}
\begin{tabular}{lrrrrr}
\toprule
code & $n$ & $k$ & $\dub$ (synd.) & $\dub$ (QDistRnd) & $kd^2/n$ \\
\midrule
\texttt{coset\_180\_20\_14\_0}        & 180 & 20 & 14 & 14 & 21.8 \\
\texttt{coset\_180\_18\_14\_1}        & 180 & 18 & 14 & 14 & 19.6 \\
\texttt{coset\_180\_16\_14\_\{2..5\}} & 180 & 16 & 14 & 14 & 17.4 \\
\bottomrule
\end{tabular}
\end{center}
Every code carries an explicit weight-$14$ witness operator, so the values are
upper bounds: the estimated parameters are $[[180,k,\le 14]]$, not certified
distances. The efficiencies $kd^2/n$ assume $d=14$ (consistent with both methods
saturating there and with the slope corroboration of Appendix~\ref{app:example});
they would beat the standard bivariate-bicycle frontier of $\approx 13.5$ if the
distance is confirmed.

\begin{thebibliography}{9}
\bibitem{Survey} \emph{Distance-Finding Algorithms for Quantum Codes and
Circuits}, arXiv:2603.22532 (2026). Taxonomy of exact (Brouwer--Zimmermann, MIP,
MaxSAT, connected-cluster) and heuristic (random information set / QDistRnd,
QDistEvol, BP\,+\,OSD syndrome decoder, Stim undetectable-error) methods.
\bibitem{QDistRnd} L.~P.~Pryadko et al., \emph{QDistRnd: A GAP package for
computing the distance of quantum error-correcting codes}, J.~Open Source Softw.\
\textbf{7}, 4120 (2022). Random-information-set upper bound.
\bibitem{RoffeBPOSD} J.~Roffe, D.~R.~White, S.~Burton, E.~Campbell,
\emph{Decoding across the quantum low-density parity-check code landscape},
Phys.\ Rev.\ Research \textbf{2}, 043423 (2020); \texttt{ldpc} package.
\bibitem{PanteleevKalachev} P.~Panteleev, G.~Kalachev,
\emph{Degenerate quantum LDPC codes with good finite length performance},
Quantum \textbf{5}, 585 (2021). Ordered-statistics decoding.
\bibitem{DKLP} E.~Dennis, A.~Kitaev, A.~Landahl, J.~Preskill,
\emph{Topological quantum memory}, J.~Math.\ Phys.\ \textbf{43}, 4452 (2002).
Sub-threshold scaling.
\bibitem{Google} Google Quantum AI, \emph{Quantum error correction below the
surface code threshold}, Nature \textbf{638}, 920 (2025); arXiv:2408.13687.
$\Lambda$-suppression and effective distance.
\bibitem{Stim} C.~Gidney, \emph{Stim: a fast stabilizer circuit simulator},
Quantum \textbf{5}, 497 (2021). Circuit-level / undetectable-error methods.
\end{thebibliography}

\end{document}
